% Options for packages loaded elsewhere
\PassOptionsToPackage{unicode}{hyperref}
\PassOptionsToPackage{hyphens}{url}
\PassOptionsToPackage{dvipsnames,svgnames,x11names}{xcolor}
%
\documentclass[
  10pt,
]{article}
\usepackage{amsmath,amssymb}
\usepackage{iftex}
\ifPDFTeX
  \usepackage[T1]{fontenc}
  \usepackage[utf8]{inputenc}
  \usepackage{textcomp} % provide euro and other symbols
\else % if luatex or xetex
  \usepackage{unicode-math} % this also loads fontspec
  \defaultfontfeatures{Scale=MatchLowercase}
  \defaultfontfeatures[\rmfamily]{Ligatures=TeX,Scale=1}
\fi
\usepackage{lmodern}
\ifPDFTeX\else
  % xetex/luatex font selection
  \setmainfont[]{DejaVu Serif}
  \setsansfont[]{DejaVu Sans}
  \setmonofont[]{DejaVu Sans Mono}
  \setmathfont[]{Latin Modern Math}
\fi
% Use upquote if available, for straight quotes in verbatim environments
\IfFileExists{upquote.sty}{\usepackage{upquote}}{}
\IfFileExists{microtype.sty}{% use microtype if available
  \usepackage[]{microtype}
  \UseMicrotypeSet[protrusion]{basicmath} % disable protrusion for tt fonts
}{}
\makeatletter
\@ifundefined{KOMAClassName}{% if non-KOMA class
  \IfFileExists{parskip.sty}{%
    \usepackage{parskip}
  }{% else
    \setlength{\parindent}{0pt}
    \setlength{\parskip}{6pt plus 2pt minus 1pt}}
}{% if KOMA class
  \KOMAoptions{parskip=half}}
\makeatother
\usepackage{xcolor}
\usepackage[margin=2.0cm]{geometry}
\usepackage{longtable,booktabs,array}
\usepackage{calc} % for calculating minipage widths
% Correct order of tables after \paragraph or \subparagraph
\usepackage{etoolbox}
\makeatletter
\patchcmd\longtable{\par}{\if@noskipsec\mbox{}\fi\par}{}{}
\makeatother
% Allow footnotes in longtable head/foot
\IfFileExists{footnotehyper.sty}{\usepackage{footnotehyper}}{\usepackage{footnote}}
\makesavenoteenv{longtable}
\setlength{\emergencystretch}{3em} % prevent overfull lines
\providecommand{\tightlist}{%
  \setlength{\itemsep}{0pt}\setlength{\parskip}{0pt}}
\setcounter{secnumdepth}{-\maxdimen} % remove section numbering
% doc_header.tex -- LaTeX preamble additions used by build_docs.sh for every
% Markdown -> LaTeX -> PDF conversion (supplements, README, submission files).
\usepackage{etoolbox}
% longtable may not be loaded by pandoc when a document has no table
\usepackage{longtable}
% bold vectors: \bm -> \symbf under unicode-math (XeLaTeX), bm package otherwise
\makeatletter
\@ifpackageloaded{unicode-math}{\providecommand{\bm}[1]{\symbf{#1}}}{\usepackage{bm}}
\makeatother
% smaller type in long tables so that wide audit tables fit the text width;
% inside a table a line break is allowed after an underscore of plain text
% (pandoc writes it as \_), so that long identifiers wrap inside their column
% rows are set a little apart, so that symbols with indices do not touch
\AtBeginEnvironment{longtable}{\footnotesize\renewcommand{\arraystretch}{1.45}%
  \let\sptunderscore\_\def\_{\sptunderscore\allowbreak}}
\setlength{\LTcapwidth}{\textwidth}
% allow line breaks after slashes and hyphens inside table cells
\usepackage{xurl}
% code blocks: break over-long lines instead of running into the margin.
% Plain ``` blocks are typeset by pandoc as {verbatim}, highlighted ones
% (```bash) as {Highlighting}; both are re-declared as breakable fancyvrb
% environments.  The continuation line is marked with a small hooked arrow.
\usepackage{fvextra}
\RecustomVerbatimEnvironment{verbatim}{Verbatim}{breaklines,breakanywhere,fontsize=\small}
\makeatletter
\@ifundefined{Highlighting}{}{%
  \RecustomVerbatimEnvironment{Highlighting}{Verbatim}{breaklines,breakanywhere,fontsize=\small,commandchars=\\\{\}}}
\makeatother
% inline code (pandoc writes it as \texttt{...}): never hyphenate -- an inserted
% hyphen could be read as part of a file name -- but allow a line break after
% an underscore, slash, comma, colon or equals sign, so that long paths and
% column lists wrap inside the text width and inside table columns.
% Inside a formula (e.g. the name of a notebook symbol in an equation) the same
% command must first leave math mode, which \text of amsmath does.
\ExplSyntaxOn
\tl_new:N \l__spt_code_tl
\cs_new_protected:Npn \__spt_texttt:n #1
  {
    \tl_set:Nn \l__spt_code_tl {#1}
    \regex_replace_all:nnN { ([\/\,\:\=]) } { \1 \c{allowbreak} } \l__spt_code_tl
    \regex_replace_all:nnN { \c{_} } { \c{_} \c{allowbreak} } \l__spt_code_tl
    \group_begin: \ttfamily \hyphenchar \font = -1 \scan_stop: \l__spt_code_tl \group_end:
  }
\RenewDocumentCommand \texttt { m }
  {
    \mode_if_math:TF { \text { \__spt_texttt:n {#1} } } { \__spt_texttt:n {#1} }
  }
\ExplSyntaxOff
% A figure is written in the Markdown source as an image paragraph followed by
% a caption paragraph ("**Figure S1.** ..."), so that the caption is also
% visible on GitHub.  An image that starts a paragraph is centered.
\usepackage{graphicx}% pandoc loads it only for documents that contain an image
\NewCommandCopy\sptincludegraphics\includegraphics
\RenewDocumentCommand\includegraphics{O{} m}{%
  \ifvmode\begin{center}\sptincludegraphics[#1]{#2}\end{center}%
  \else\sptincludegraphics[#1]{#2}\fi}
% keep figures where they are referenced
\usepackage{float}
\floatplacement{figure}{H}
% running header with the document title
\usepackage{fancyhdr}
\pagestyle{fancy}
\fancyhf{}
\fancyhead[L]{\small PROJECT spacetimes --- documentation and provenance}
\fancyhead[R]{\small\thepage}
\renewcommand{\headrulewidth}{0.4pt}
% 16x16 matrices: tighter columns so that a full matrix fits the text width
\setlength{\arraycolsep}{3.2pt}
% The documents of this project contain long tagged symbols, file paths and
% identifiers that cannot be hyphenated.  A paragraph is therefore allowed to
% have looser lines before TeX gives up and lets a line stick out of the margin.
\AtBeginDocument{\tolerance=2000 \emergencystretch=5em \hbadness=2000}
\ifLuaTeX
  \usepackage{selnolig}  % disable illegal ligatures
\fi
\IfFileExists{bookmark.sty}{\usepackage{bookmark}}{\usepackage{hyperref}}
\IfFileExists{xurl.sty}{\usepackage{xurl}}{} % add URL line breaks if available
\urlstyle{same}
\hypersetup{
  pdftitle={PROJECT spacetimes --- folder Revision\_2b: right-handed Dirac spinors},
  pdfauthor={Prepared for Patrick Lee Nash, Ph.D.},
  colorlinks=true,
  linkcolor={blue},
  filecolor={Maroon},
  citecolor={blue},
  urlcolor={blue},
  pdfcreator={LaTeX via pandoc}}

\title{PROJECT spacetimes --- folder Revision\_2b: right-handed Dirac
spinors}
\usepackage{etoolbox}
\makeatletter
\providecommand{\subtitle}[1]{% add subtitle to \maketitle
  \apptocmd{\@title}{\par {\large #1 \par}}{}{}
}
\makeatother
\subtitle{What the folder contains, its state, and how to rebuild it}
\author{Prepared for Patrick Lee Nash, Ph.D.}
\date{November 15, 2026}

\begin{document}
\maketitle

\hypertarget{what-the-folder-is-for}{%
\section{1. What the folder is for}\label{what-the-folder-is-for}}

The author instructed, on November 15, 2026, that every instruction,
calculation, analysis, document, manuscript and provenance of the
folders \texttt{Revision\_Deflation\_CJT\_and\_Kohn-Sham} and
\texttt{Revision\_Deflation\_CJT\_and\_Kohn-Sham\_definition\_C} be
repeated in this folder for right-handed Dirac spinors, that is, for a
field \(\psi\) with \(P^R\psi=\psi\), where \(P^R=\tfrac12(1+\chi)\) and
\(\chi\) is the chirality matrix
(\(\chi=\gamma^{0}\gamma^{1}\cdots\gamma^{7}\) on the 4+4 spacetime,
\(\chi=i\gamma^{0}\gamma^{1}\cdots\gamma^{7}\) on the 7+1 spacetime,
with the Dirac matrices \(\gamma^A\) of the project).

\hypertarget{the-state-of-the-folder-paused}{%
\section{2. The state of the folder:
paused}\label{the-state-of-the-folder-paused}}

Before anything was repeated, it was computed, in exact arithmetic,
which terms of the model of the two folders a right-handed field keeps.
On both spacetimes it keeps the kinetic term and half of the invariants
of the self-interaction, but \textbf{not the mass term, not the coupling
to the Vorton-Scalar, and not the pair along the hidden axis \(x0\)}
whose condensate is the singlet superfluid of the two folders (29
checks, all passing). The instruction can therefore not be carried out
as a repetition. As the author's Hard Rule 8 requires, the provenance
\texttt{P0} explains why, says what the result implies for each
calculation of the two folders, and states the decision that is the
author's. The work is paused until he decides.

\hypertarget{contents}{%
\section{3. Contents}\label{contents}}

\begin{longtable}[]{@{}
  >{\raggedright\arraybackslash}p{(\columnwidth - 2\tabcolsep) * \real{0.4337}}
  >{\raggedright\arraybackslash}p{(\columnwidth - 2\tabcolsep) * \real{0.5663}}@{}}
\toprule\noalign{}
\begin{minipage}[b]{\linewidth}\raggedright
path
\end{minipage} & \begin{minipage}[b]{\linewidth}\raggedright
content
\end{minipage} \\
\midrule\noalign{}
\endhead
\bottomrule\noalign{}
\endlastfoot
\texttt{provenance/P0\_right\_handed\_pause} (\texttt{.md},
\texttt{.tex}, \texttt{.pdf}) & the instruction; what ``right-handed''
means; the calculation and all its results; what they imply for each
command of the two folders; why the work is paused and the author's
decision; notation \\
\texttt{README} (\texttt{.md}, \texttt{.tex}, \texttt{.pdf}) & this
list \\
\texttt{code/chirality-rs} & the Rust program \texttt{chirality}: the
exact tests of every term of the model for both chiralities and both
spacetimes \\
\texttt{code/scripts/fill\_docs.py} & writes the Markdown documents from
their sources in \texttt{docs\_src}, filling in the results from
\texttt{data} \\
\texttt{data/chirality\_checks.txt} & every check of the program with
its statement and verdict \\
\texttt{data/chirality\_terms.csv} & every term of the model: kept or
lost, for each spacetime and chirality \\
\texttt{data/chirality\_bilinears\_by\_rank.csv} & the bilinears of the
self-interaction that are kept, by rank \\
\texttt{data/chirality\_pairs\_by\_rank.csv} & the pair operators that
are kept, by rank \\
\texttt{docs\_src} & the sources of the two documents \\
\end{longtable}

\hypertarget{how-to-rebuild-everything}{%
\section{4. How to rebuild everything}\label{how-to-rebuild-everything}}

From the folder \texttt{Revision\_2b} (a Rust toolchain is needed; no
crate is downloaded; the program uses the library of the main project,
\texttt{code/spacetimes-rs} at the root of the repository):

\begin{verbatim}
cargo run --release --manifest-path code/chirality-rs/Cargo.toml
python3 code/scripts/fill_docs.py
\end{verbatim}

Then, from the root of the repository (pandoc, XeLaTeX and latexmk are
needed):

\begin{verbatim}
bash code/scripts/build_docs.sh Revision_2b/provenance/P0_right_handed_pause.md \
     Revision_2b/README.md
\end{verbatim}

The program ends with ``chirality: N of N checks passed'' and returns
the exit status 1 if a check fails.

\end{document}
